\chapter{Advection–diffusion: when Galerkin fails, and SUPG}
\label{ch:C2_advection_diffusion}
\noindent\texttt{tutorials/C\_time/C2\_advection\_diffusion.py}\\[0.75em]
\begin{outcome}
You can recognize the failure mode of plain Galerkin on
advection-dominated transport (node-to-node oscillations governed by the
cell Peclet number), fix it with SUPG — a residual-based stabilization
that adds diffusion only ALONG the streamline and only where the residual
is nonzero — and verify that the fix (i) kills the wiggles and (ii) does
NOT degrade the MMS convergence order (that is what "consistent
stabilization" means). p1/p2 in space, BDF1/BDF2 in time.
\end{outcome}

u\_t + a.grad(u) - kappa lap(u) = f. The SUPG-stabilized weak
form adds, per element,
    sum\_e ( tau (a.grad w),  R(u) )\_e,   R = u\_t + a.grad u - kappa lap u - f
with tau from the same VMS toolbox track D uses (diffsim.physics.vms,
metric form). At p1 the lap(u\_h) term in R vanishes elementwise (Q1's
diagonal second derivatives are identically zero — the same fact behind
the m1a shift-order rule). Steady MMS first, then transient with BDF.

\begin{expected}
\begin{verbatim}
all MEASURED, and each number teaches a theorem):
    steady, Galerkin:  overshoot +3.2e-2 above the exact max (wiggles)
    steady, SUPG p1:   overshoot +1.3e-3; orders 1.73 -> 1.53 — NOT 2, and
        correctly so: the classical SUPG estimate in the advection-dominated
        limit is O(h^{k+1/2}) — order 1.5 for p1. You are measuring a
        theorem, not a bug.
    steady, SUPG p2:   orders ~1 — DEGRADED, because this chapter's residual
        drops the kappa*lap(u_h) term, which is legal at p1 (Q1 diagonal
        second derivatives vanish) but INCONSISTENT at p2. Fixing it is
        Explore (b) — the measurement is left broken on purpose.
    transient + SUPG (p1, dt-independent tau, per-order LU — read the
    care point in march()): BDF1 rates 0.97/0.99, BDF2 rates 2.02/2.07
\end{verbatim}
\end{expected}

\section*{The script}
\begin{lstlisting}
import numpy as np
import scipy.sparse as sp
from scipy.sparse.linalg import splu

from diffsim.octree.build import build_uniform
from diffsim.mesh.nodes import build_mesh
from diffsim.mesh.constraints import build_constraints
from diffsim.mesh.basis import basis_tables
from diffsim.assembly.operators import DeviceMesh, assemble_csr
from diffsim.physics.poisson import gauss_points, l2_error_masked
from diffsim.physics.vms import tau_metric_host
from diffsim.solvers.timestepping import bdf_coeffs, bdf_order_now, History

DEVICE = "cuda:0"
KAPPA = 1e-3
AVEC = np.array([1.0, 0.5]) / np.sqrt(1.25)      # unit advection direction


class AdvDiff:
    """Host-assembled SUPG advection-diffusion on a uniform 2-D mesh.
    (Track D does this fused in warp kernels; here the einsums ARE the
    lesson — read them against the weak form.)"""

    def __init__(self, level, p, dt=None, order=2):
        tree = build_uniform(level, dim=2)
        self.mesh = build_mesh(tree, p=p)
        self.cons = build_constraints(self.mesh)
        self.dm = DeviceMesh.from_mesh(self.mesh, self.cons,
                                       basis_tables(p, dim=2), DEVICE)
        self.T = self.cons.T.tocsr()
        self.K = assemble_csr(self.dm) * KAPPA     # kappa (grad w, grad u)
        self.xq = gauss_points(self.mesh, self.dm.tables_by_p)
        self.coords = self.mesh.node_coords[self.cons.free_nodes]
        self.bdry = np.where(
            self.mesh.boundary_nodes[self.cons.free_nodes])[0]
        self.dt, self.order = dt, order
        sigma = (bdf_coeffs(order, dt)[0] / dt) if dt else 0.0
        self._assemble_terms(sigma)

    def _assemble_terms(self, sigma):
        """conv + SUPG matrices and the SUPG-weighted test operator for the
        RHS: all einsum, all commented."""
        dm, mesh = self.dm, self.mesh
        Nn = dm.n_nodes
        C = sp.lil_matrix((Nn, Nn))
        Msupg = sp.lil_matrix((Nn, Nn))            # tau (a.grad w, sigma u)
        Mass = sp.lil_matrix((Nn, Nn))
        rows, cols = [], []
        Cv, Sv, Mv = [], [], []
        self._supg_test = {}                       # per-bin (a.grad w) tau
        for pv in dm.bins:
            tb = dm.tables_by_p[pv]
            h = mesh.tree.h()[mesh.bins[pv]]
            jac = (h / 2.0) ** dm.dim
            dsc = 2.0 / h
            conn = mesh.conn_of[pv].astype(np.int64)
            nbf = conn.shape[1]
            # physical gradients: dN[q,a,d] * dscale_e  -> per-element
            aGN = np.einsum("qad,d->qa", tb.dN, AVEC)      # a.grad N (ref)
            tau, _ = tau_metric_host(1.0, h, KAPPA, dt=None, dim=2,
                                     timestab=False)       # per element
            # NOTE: tau is kept dt-INDEPENDENT even in transient runs. With
            # the transient term in tau, every dt in a convergence ladder
            # solves a slightly different spatially-stabilized problem and
            # the ladder floors (measured here: BDF rates ~0; same effect as
            # m1b findings 1). Production exposes this as the timeStab
            # toggle; ladders run it off.
            # (w, a.grad u):  N_a * aGN_b * dsc
            Ce = np.einsum("qa,qb,q,e->eab", tb.N, aGN, tb.w,
                           jac * dsc)
            # SUPG (tau a.grad w, a.grad u):  aGN_a aGN_b dsc^2
            Se = np.einsum("qa,qb,q,e->eab", aGN, aGN, tb.w,
                           jac * dsc ** 2 * tau)
            # SUPG mass part (tau a.grad w, sigma N_b) + Galerkin mass
            Sm = np.einsum("qa,qb,q,e->eab", aGN, tb.N, tb.w,
                           jac * dsc * tau)
            Me = np.einsum("qa,qb,q,e->eab", tb.N, tb.N, tb.w, jac)
            rows.append(np.repeat(conn, nbf, axis=1).ravel())
            cols.append(np.tile(conn, (1, nbf)).ravel())
            Cv.append(Ce.ravel())
            Sv.append(Se.ravel())
            # the SUPG-weighted mass (tau a.grad w, N_b) must multiply the
            # WHOLE discrete time derivative — sigma u^{n+1} on the LHS AND
            # the history term on the RHS. Folding it into the LHS only is
            # an O(tau) inconsistency (measured: erratic dt-ladders).
            Mv.append((Me + Sm).ravel())
            # RHS test weights: (N_a + tau a.grad N_a dsc) per GP
            self._supg_test[pv] = (tb.N[None, :, :]
                                   + (tau * dsc)[:, None, None]
                                   * aGN[None, :, :])       # [ne,q,a]
        mk = lambda vals: sp.coo_matrix(
            (np.concatenate(vals),
             (np.concatenate(rows), np.concatenate(cols))),
            shape=(Nn, Nn)).tocsr()
        self.C = self.T.T @ mk(Cv) @ self.T
        self.S = self.T.T @ mk(Sv) @ self.T
        self.M = self.T.T @ mk(Mv) @ self.T   # Galerkin + SUPG-test mass

    def rhs(self, f_gp_by_bin):
        """(w + tau a.grad w, f) — the SUPG-consistent load."""
        F = np.zeros(self.dm.n_nodes)
        for pv in self.dm.bins:
            h = self.mesh.tree.h()[self.mesh.bins[pv]]
            jac = (h / 2.0) ** self.dm.dim
            tb = self.dm.tables_by_p[pv]
            ne = len(h)
            fq = f_gp_by_bin[pv].reshape(ne, tb.nqp)
            be = np.einsum("eqa,eq,q,e->ea", self._supg_test[pv], fq,
                           tb.w, jac)
            np.add.at(F, self.mesh.conn_of[pv].ravel(), be.ravel())
        return np.asarray(self.T.T @ F)

    def system(self, supg=True):
        A = self.K + self.C + (self.S if supg else 0.0)
        if self.dt:
            A = A + (bdf_coeffs(self.order, self.dt)[0] / self.dt) * self.M
        return A


def steady_case(level, p, supg):
    us = lambda x: np.sin(np.pi * x[:, 0]) * np.sin(np.pi * x[:, 1])

    def grad_us(x):
        return np.stack([np.pi * np.cos(np.pi * x[:, 0])
                         * np.sin(np.pi * x[:, 1]),
                         np.pi * np.sin(np.pi * x[:, 0])
                         * np.cos(np.pi * x[:, 1])], axis=1)
    fs = lambda x: (grad_us(x) @ AVEC) + 2 * np.pi ** 2 * KAPPA * us(x)
    ad = AdvDiff(level, p)
    A = ad.system(supg=supg).tolil()
    b = ad.rhs({pv: fs(ad.xq[pv]) for pv in ad.xq})
    for i in ad.bdry:
        A.rows[i] = [int(i)]
        A.data[i] = [1.0]
        b[i] = 0.0
    u = splu(A.tocsr().tocsc()).solve(b)
    u_all = np.asarray(ad.T @ u)
    err = l2_error_masked(ad.dm, u_all, us, lambda x: np.ones(len(x), bool))
    return err, float(u.max() - 1.0)             # overshoot above exact max


if __name__ == "__main__":
    e_g, over_g = steady_case(5, 1, supg=False)
    e_s, over_s = steady_case(5, 1, supg=True)
    print(f"steady level 5 p1:  Galerkin overshoot {over_g:+.3e}   "
          f"SUPG overshoot {over_s:+.3e}")
    for p, levels in ((1, (4, 5, 6)), (2, (3, 4, 5))):
        errs = [steady_case(lv, p, True)[0] for lv in levels]
        orders = [np.log2(errs[i] / errs[i + 1]) for i in range(2)]
        print(f"SUPG steady p={p}: errors "
              + "  ".join(f"{e:.3e}" for e in errs)
              + "   orders " + "  ".join(f"{o:.2f}" for o in orders))

    # transient: u* = cos(2 pi t) sin(pi x) sin(pi y), same-mesh reference
    OM = 2 * np.pi
    us = lambda x, t: np.cos(OM * t) * np.sin(np.pi * x[:, 0]) \
        * np.sin(np.pi * x[:, 1])

    def f_t(x, t):
        c, s_ = np.cos(OM * t), np.sin(OM * t)
        sp_ = np.sin(np.pi * x[:, 0]) * np.sin(np.pi * x[:, 1])
        gx = np.pi * np.cos(np.pi * x[:, 0]) * np.sin(np.pi * x[:, 1])
        gy = np.pi * np.sin(np.pi * x[:, 0]) * np.cos(np.pi * x[:, 1])
        return (-OM * s_ * sp_ + c * (AVEC[0] * gx + AVEC[1] * gy)
                + 2 * np.pi ** 2 * KAPPA * c * sp_)

    def march(n, order):
        T_f = 0.5
        dt = T_f / n
        ad = AdvDiff(5, 1, dt=dt, order=order)
        # CARE POINT (measured the hard way): the LHS's b0/dt MUST match
        # the RHS's BDF coefficients at every step. The bootstrap step runs
        # BDF1 even in an order-2 march, so cache one LU per EFFECTIVE
        # order — prebuilding a single BDF2 matrix injects an O(1) error
        # at step 1 that a weakly-damped problem never forgets (we measured
        # a dt-independent 0.17 'floor' before fixing this).
        lus = {}
        hist = History()
        hist.rotate(us(ad.coords, 0.0))
        t = 0.0
        for k in range(n):
            t += dt
            o = bdf_order_now(t, dt, order, have_history=hist.have(2))
            b0, b1, b2 = bdf_coeffs(o, dt)
            if o not in lus:
                A = (ad.K + ad.C + ad.S + (b0 / dt) * ad.M).tolil()
                for i in ad.bdry:
                    A.rows[i] = [int(i)]
                    A.data[i] = [1.0]
                lus[o] = splu(A.tocsr().tocsc())
            b = ad.rhs({pv: f_t(ad.xq[pv], t) for pv in ad.xq}) \
                - ad.M @ hist.rhs_history(b1, b2, dt)
            b[ad.bdry] = 0.0
            hist.rotate(lus[o].solve(b), dt=dt)
        return hist.pre1

    ref = march(512, 2)
    for order in (1, 2):
        errs = [np.sqrt(((march(n, order) - ref) ** 2).mean())
                for n in (32, 64, 128)]
        rates = [np.log2(errs[i] / errs[i + 1]) for i in range(2)]
        print(f"transient BDF{order}: errors "
              + "  ".join(f"{e:.3e}" for e in errs)
              + "   rates " + "  ".join(f"{r:.2f}" for r in rates))
\end{lstlisting}

\begin{explore}
\begin{verbatim}
(a) Sweep kappa from 1e-1 to 1e-5 at level 5 and plot the Galerkin
      overshoot vs the cell Peclet number Pe = |a| h / (2 kappa). Where is
      the classical Pe = 1 threshold in your data?
  (b) SUPG at p2 uses the same tau. The lap(u_h) residual term is NOT zero
      at p2 (mixed second derivatives exist) and we dropped it — measure
      whether that costs the third order (compare with the p2 ladder
      above). This is an open sharp edge in most production codes too.
  (c) Crank up |a| by 10x and keep tau's formula. Wiggles again? tau's
      scaling with |a| (read tau_metric_host) should protect you —
      verify.
  (d) PERFORMANCE CORNER: transient runs cache ONE LU because tau and a
      are constant. In track D the advecting field changes every step —
      re-time this loop with a per-step refactorization and estimate what
      fraction of a Navier-Stokes step budget the factorization would eat
      at level 6. (That is why D-track uses iterative solvers + the fused
      Krylov of P2.)
\end{verbatim}
\end{explore}
